\clearpage
\begin{centering}
\textbf{ABSTRACT}\\
\vspace{\baselineskip}
\end{centering}

Gradient Ascent Pulse Engineering (GRAPE) is a popular technique in quantum optimal control, and can be combined
with automatic differentiation (AD) to facilitate on-the-fly evaluation of cost-function gradients. I illustrate that the
convenience of AD comes at a significant memory cost due to the cumulative storage of a large number of states and
propagators. For quantum systems of increasing Hilbert space size, this imposes a significant bottleneck. I revisit the
strategy of hard-coding gradients in a scheme that fully avoids propagator storage and significantly reduces memory
requirements. 
I benchmark runtime and memory usage and compare this approach to AD-based implementations, with a focus on
pushing towards larger Hilbert space sizes. The results confirm that the AD-free approach facilitates the application of optimal control for large quantum systems which would otherwise be difficult to tackle. To facilitate the choice of optimal numerical strategy for GRAPE-based quantum optimal control, I formulate a decision tree. Based on the guidance from the decision tree, I develop a GRAPE-based optimization procedure for transmon. The optimizer generates detuning-robust pulses, which are crucial in the existing superconducting quantum processor unit architecture. The generated pulses outperform conventionally used transmon pulses in the aspect of infidelity robustness. 
